\section{A compiler grammar for matchertext}
\label{app:grammar}

Realizing the syntactic extensions of \cref{sec:host} in an existing compiler
reduces to a single requirement:
the additions must form a \emph{conservative extension} of the language,
under which every previously valid program stays valid and keeps its meaning.
Only a conservative extension can be adopted incrementally,
enabled behind a configuration flag without endangering existing code,
the backward compatibility that \cref{sec:host} identifies as essential to adoption.
The general technique is to guard each new form
with a marker the old grammar could never accept,
so the extension fires only on input the language previously rejected.
We give grammar rules of this kind for the two directions in turn,
first the embedded matchertext literal itself
and then the accommodation a host language must make to admit it.

\subsection{Embedded-matchertext grammar}
\label{app:grammar:embed}

We now cast the matchertext literal of \cref{sec:host:c:literals}
as concrete grammar rules that a C-like language could fold into its lexer.

We assume the target language already parses ordinary string literals
with a rule of the general shape:
\[
\begin{array}{rcl}
\mathit{content}
  & ::= & \texttt{...} \\

\mathit{delimiter}
  & ::= & \texttt{"} \\

\mathit{string}
  & ::= & \mathit{delimiter}\ \mathit{content}\ \mathit{delimiter}
\end{array}
\]
The delimiter is shown here as a double quote,
but any string delimiter the language already uses would serve equally well,
for instance \verb|'|, \verb|"|, \verb|"""|, \verb|'''|, or a backtick.

The proposed form is \verb|M|$d$\verb|"|$o\dots c$\verb|"|,
an ordinary string literal carrying a mandatory \verb|M| prefix
and an optional specifier $d$.
Its content $o\dots c$ is itself wrapped in a matcher pair,
$o$ one of \verb|(|, \verb|[|, \verb|{| and $c$ its matching close.
That wrapping is the crux:
the processor finds the literal's end by matcher balance,
at the $c$ closing the opening $o$ (\cref{thm:design:findend}),
not at the first quote it meets inside,
so the content may hold any character, quotes included,
as long as its matchers balance.
This is exactly the matchertext discipline,
and \cref{thm:design:embed} then guarantees the content embeds verbatim
without escaping.
The optional $d$, one or more strictly nested matcher pairs,
names the pair reserved for interpolation holes or antiquotations
inside the literal (\cref{app:interp}),
so a template can mark its holes unambiguously.
The \verb|M| prefix keeps the extension conservative
in languages whose grammars reject an identifier directly before a quote,
such as Python, Go, and Java,
and Rust~2021, which reserves all such prefixes for future use.
Other languages need more care.
In C and C++, a macro named \verb|M| can make \verb|M"|$\dots$\verb|"| valid today,
the same incompatibility that C++11 accepted for raw-string prefixes such as \verb|R|.
Scala reads any identifier before a quote as a string interpolator,
and JavaScript reads an identifier before a backtick as a tagged template.
Such languages need a marker that their grammar rejects,
as CHICKEN Scheme does with \verb|#M| (\cref{sec:impl}).

We first define the matchertext symbols used (prefixed with MT):
\[
\begin{array}{rcl}
\mathit{MT\_prefix}
  & ::= & \texttt{M} \\

\mathit{MT\_d}
  & ::= & \texttt{(}\mathit{MT\_d}\texttt{)}
    \mid \texttt{[}\mathit{MT\_d}\texttt{]}
    \mid \texttt{\{}\mathit{MT\_d}\texttt{\}} \mid \epsilon \\

\mathit{MT\_string}
  & ::= & \mathit{delimiter}\ \texttt{(}\,\mathit{MT\_body}\,\texttt{)}\ \mathit{delimiter} \\
  & \mid & \mathit{delimiter}\ \texttt{\{}\,\mathit{MT\_body}\,\texttt{\}}\ \mathit{delimiter} \\
  & \mid & \mathit{delimiter}\ \texttt{[}\,\mathit{MT\_body}\,\texttt{]}\ \mathit{delimiter} \\

\mathit{MT\_body}
  & ::= & \mathit{MT\_char}\ \mathit{MT\_body}
    \mid \texttt{(}\,\mathit{MT\_body}\,\texttt{)}\,\mathit{MT\_body} \\
  & \mid & \texttt{[}\,\mathit{MT\_body}\,\texttt{]}\,\mathit{MT\_body}
    \mid \texttt{\{}\,\mathit{MT\_body}\,\texttt{\}}\,\mathit{MT\_body}
    \mid \epsilon \\

\mathit{MT\_char}
  & ::= & \textrm{any character except } \texttt{(}\ \texttt{)}\ \texttt{[}\ \texttt{]}\ \texttt{\{}\ \texttt{\}} \\
\end{array}
\]
\verb|MT_body| is the matchertext language $L$ of \cref{sec:design:abstract}
in the standard configuration.
It replaces the ordinary \verb|content|,
so quotes and escapes inside it have no meaning to the host.

We then introduce the extended rule \verb|string_ext|
and use it wherever the grammar previously accepted a \verb|string|.
Its first alternative is the unchanged \verb|string| rule,
so every existing literal parses exactly as before,
and its remaining alternatives add the matchertext forms:
\[
\begin{array}{rcl}
\mathit{string\_ext}
  & ::= & \mathit{string} \\
  & \mid & \mathit{MT\_prefix}\,\mathit{MT\_d}\,\mathit{MT\_string}
\end{array}
\]

We also give the same rules in LL(1) form,
so a predictive parser can select each alternative from one token of lookahead.
Here \verb|MT_string| omits the leading delimiter, which \verb|MT_content| supplies,
letting the parser decide as soon as it sees the opening quote or an opening
matcher after the \verb|M|:
\[
\begin{array}{rcl}
\mathit{string\_ext}
  & ::= & \mathit{string} \mid \mathit{MT\_prefix}\ \mathit{MT\_content} \\

\mathit{MT\_content}
  & ::= & \mathit{MT\_d}\ \mathit{delimiter}\ \mathit{MT\_string} \\

\mathit{MT\_d}
  & ::= & \texttt{(}\ \mathit{MT\_d}\ \texttt{)}
          \mid \texttt{[}\ \mathit{MT\_d}\ \texttt{]} \\
  & \mid & \texttt{\{}\ \mathit{MT\_d}\ \texttt{\}} \mid \epsilon \\

\mathit{MT\_string}
  & ::= & \texttt{(}\,\mathit{MT\_body}\,\texttt{)}\ \mathit{delimiter} \\
  & \mid & \texttt{\{}\,\mathit{MT\_body}\,\texttt{\}}\ \mathit{delimiter} \\
  & \mid & \texttt{[}\,\mathit{MT\_body}\,\texttt{]}\ \mathit{delimiter} \\
\end{array}
\]

\subsection{Host-language grammar}
\label{app:grammar:host}

Hosting embedded matchertext asks less of the grammar than embedding does,
at least for C-like languages,
because the property a host must enforce inside an embedding,
that matchers balance,
is one these languages already enforce throughout their ordinary syntax.
The expression, block, and argument-list grammars of
C, C++, Rust, Java, Python, and their relatives
all require every \verb|(|, \verb|[|, and \verb|{| to be matched
for a program to compile at all,
so at the structural level their code is already matchertext.
The only contexts in which they admit unmatched matchers today
are the opaque ones, string and character literals and comments,
whose contents the compiler does not otherwise interpret.

Those contexts are also the ones backward compatibility forbids us
to constrain for existing, non-embedded code,
so a host does not tighten them globally.
It instead introduces a bracketed embedding region within them,
such as the \verb|\[|$m$\verb|]| escape of \cref{sec:host:c:escapes},
inside which two rules apply,
following the general hosting principle of \cref{sec:host}:
the host disables its own escaping and other transformations,
and locates the region's end by matcher balance (\cref{thm:design:findend})
rather than by its usual literal or comment terminator.
Because the surrounding code already balances its matchers,
this region composes cleanly with it,
and the grammar change a C-like host needs is correspondingly small:
the discipline it must impose inside an embedding
is the one it already imposes everywhere else.